\subsection{All torsion levels for tame potentially good elliptic curves}

\begin{theorem}[B-RD12: the sharp tame higher-level criterion]\label{thm:B-RD12}
\status{Proved; independently audited by the parent agent.}
Fix $p\geq5$, $n\geq1$, and $m\in\{2,3,4,6\}$ with $m\mid p+1$.
Assume $e$ is a positive multiple of $p-1$, and put
$s=v_p(e/(p-1))$. There is a complete discretely valued $K$ of absolute
ramification index $e$ and an elliptic curve with full rational
$p^n$-torsion and semistability defect $m$ if and only if
\[
 s\geq n-1,\qquad
 e\geq\frac{p^2-1}{m}\,
        p^{\max(2n-2,\,4n-4-s)}.                  \tag{B.12.1}
\]
Arbitrary perfect residue fields are allowed; the constructions use
$\overline{\mathbf F}_p$. The same assertion holds for $p=3$ with a
specified tame defect $m=2$ or $4$. In that assertion no restriction on
wild defects is claimed here; their sharp higher-level criterion is proved separately in Theorem~\ref{thm:B-RD26-wild}.
\end{theorem}
\begin{proof}
We first obtain an exact integer criterion. Pass to completed maximal
unramified bases and a tame cyclic degree-$m$ good-reduction extension
$L/K$. Use the eigenparameter from Theorem~\ref{thm:B-RD7}.
Let $A$ and $B$ be the maximum and minimum valuations of a primitive
$p^n$-torsion point. Points $P,Q$ achieving them form a basis if $A>B$:
otherwise $Q=aP+pR$ with $a$ a unit; multiplication by $p$ strictly
increases every positive formal valuation, contradicting minimality
of $B$. If $A=B$, choose any basis. We have
\[
 A\geq B\geq1,\qquad A\equiv B\equiv q\pmod m,
 \qquad \gcd(q,m)=1.                              \tag{B.12.2}
\]

The full-level deformation ring has the factorization of
\cite[Proposition~2.2.3(vii)]{CarterStrauch2021LubinTate}.
Iterated $j$ times, it expresses the coordinate of $[p^j]Q$ as a unit
times a product of $p^j$ primitive level-$n$ coordinates, each raised
to $p^j$. The relevant lines have representatives $(a,1)$ with
$a\equiv0\pmod {p^{n-j}}$. All these coordinates have valuation $B$,
because the contribution $[a]P$ has valuation strictly greater than
$B$. Consequently
\[
 v_L(t([p^j]Q))=p^{2j}B\quad(0\leq j\leq n-1).   \tag{B.12.3}
\]
This factorization applies to the actual good elliptic curve by its
map from the universal framed deformation ring with full level.

Primitive vectors modulo scalar units are the lines of
$\mathbf P^1(\mathbf Z/p^n)$. There are $p^n$ representatives $(a,1)$,
all of valuation $B$: when $A=B$, a cancellation would contradict the
chosen maximum $A$. There is the line $(1,0)$, of valuation $A$.
For $1\leq j<n$, there are $(p-1)p^{n-j-1}$ representatives $(1,b)$
with $v_p(b)=j$, each of valuation $\min(A,p^{2j}B)$. Again equality
cannot cause cancellation beyond $A$, since $A$ is the maximum among
primitive points. The series $[p^n](T)/[p^{n-1}](T)$ has constant term
$p$, distinguished degree $(p^2-1)p^{2n-2}$, and precisely the
primitive $p^n$-torsion roots. Taking its constant term gives
\[
 me=(p-1)p^{n-1}\left(
 A+p^nB+\sum_{j=1}^{n-1}(p-1)p^{n-j-1}
                \min(A,p^{2j}B)\right).          \tag{B.12.4}
\]

We prove that every pair of integers satisfying (B.12.2) and (B.12.4)
can be realized; for the stated range it will suffice to use $B=q=1$.
Use the full Drinfeld $p^n$-level deformation ring $R_n$, whose regular
parameters $x,y$ are the universal basis coordinates
\cite[Theorem~2.1.2 and Proposition~2.2.3(i)]{CarterStrauch2021LubinTate}.
The eigenparameter averaging and equivariant Cohen presentation of
Theorem~\ref{thm:B-RD7} apply verbatim to this level.
We justify the weighted specialization, including its possible
cancellations, as follows.

Work modulo $p$ in the power-series presentation and give $X,Y$
weights $A,B$. Lift the full-level factorization (B.12.3), rather than
an arbitrary expression for $[p^j]y$. Each factor with representative
$(a,1)$ and $a\equiv0\pmod {p^{n-j}}$ has weighted initial form $Y$,
up to a nonzero constant. Therefore a lift of $[p^j]y$ has weighted
initial form $c_jY^{p^{2j}}$, with $c_j\in k^\times$. It follows that
for a primitive line $(1,b)$, $v_p(b)=j$, one can choose a lift whose
initial form is the lower-weight term of
\[
 X+c_bY^{p^{2j}},\qquad c_b\ne0,
\]
or that entire binomial when the two weights agree. The other
primitive lines are handled by their linear terms. These choices are
legitimate lifts because the factorization identities hold in $R_n$;
changes by the defining relation do not alter their images.

The constant-term identity for $[p^n](T)/[p^{n-1}](T)$ gives a generator
of the Cohen kernel of the form $p$ minus a unit times the product of
all primitive coordinates. Average this generator under $C_m$.
Its weighted initial form modulo $p$ is preserved: every initial
factor has scalar character $\zeta$, since
$p^{2j}\equiv1\pmod m$, and the total number
$(p^2-1)p^{2n-2}$ of factors is divisible by $m$.
Now specialize $X=\pi^A$, $Y=c\pi^B$, choosing the residue of $c$
to avoid the finitely many roots of the displayed initial forms.
The resulting relation has linear term $p$ and first nonzero term
modulo $p$ of degree $d$, where $d$ is the right side of (B.12.4).
It defines a complete DVR with $v_L(p)=d$.
Every primitive-coordinate weight is strictly smaller than $d$, so
terms in its chosen lift divisible by $p$ cannot change that weight.
Thus all the asserted valuations hold after specialization, with no
circular assumption on $v_L(p)$.

For $A\equiv B\equiv1\pmod m$, the action $\pi\mapsto\zeta\pi$
preserves the relation and the specialization. Tame descent gives a
curve over the invariant field, of absolute ramification $d/m$,
with full rational $p^n$-torsion and exact defect $m$, as in
Theorem~\ref{thm:B-RD7}. For a general common residue $q$ in
(B.12.2), replace the action on $\pi$ by multiplication by
$\zeta^{q^{-1}}$, where the inverse is taken modulo $m$.

It remains to simplify (B.12.4), which also proves that $B=1$ suffices.
Choose $r\in\{0,\ldots,n-1\}$ so that
$p^{2r}B\leq A\leq p^{2r+2}B$, with the upper bound omitted for
$r=n-1$. Summing the geometric series in (B.12.4) gives
\[
 me=(p-1)p^{2n-r-2}(A+p^{2r+1}B).                \tag{B.12.5}
\]
It implies the two conditions
\[
 p^{2n-r-2}(p-1)\mid e,\qquad
 e\geq \frac{p^2-1}{m}p^{2n+r-2}.                \tag{B.12.6}
\]
Conversely suppose these conditions hold for some $r$. Increase $r$
until it is $n-1$, or until the second inequality fails for $r+1$;
the divisibility remains true. Put
\[
 B=1,\qquad
 A=\frac{me}{(p-1)p^{2n-r-2}}-p^{2r+1}.
\]
Then $A$ is an integer, $A\equiv1\pmod m$, and
$p^{2r}\leq A<p^{2r+2}$ unless $r=n-1$, when only the lower bound
is needed. Hence it realizes (B.12.4) by the construction above.
Finally the least permissible $r$ is
$\max(0,2n-2-s)$; such an $r\leq n-1$ exists exactly when
$s\geq n-1$. Substitution gives (B.12.1).
\end{proof}

\begin{corollary}[B-RD13: the least good-forcing level in the tame case]
\label{cor:B-RD13}
\status{Proved; independently audited by the parent agent.}
Let $p\geq5$, $e\in(p-1)\mathbf Z_{>0}$, and let $m_{\max}(p)$ be
as in Theorem~\ref{thm:B-RD7}. If $m_{\max}(p)e<p^2-1$, put $N=0$.
Otherwise put
\[
 s=v_p(e/(p-1)),\qquad
 L=\left\lfloor\log_p\frac{m_{\max}(p)e}{p^2-1}\right\rfloor,
 \qquad
 N=1+\min\left(s,\left\lfloor\frac L2\right\rfloor,
                     \left\lfloor\frac{L+s}{4}\right\rfloor\right).
\]
Among potentially good elliptic curves over all such fields of
absolute ramification $e$, the least positive level $n$ at which
full rational $p^n$-torsion forces good reduction is $N+1$.
For $p=3$, the same formula with $m_{\max}=4$ gives the least level
within the class of potentially good curves with tame inertia image.
\end{corollary}
\begin{proof}
For a fixed level, the largest allowed defect gives the weakest
inequality in Theorem~\ref{thm:B-RD12}. Its inequalities are equivalent
to $n-1\leq s$, $2n-2\leq L$, and $4n-4-s\leq L$.
Solving them gives exactly the displayed largest bad level $N$.
The construction realizes it, and its lower levels, while every
potentially good curve outside the bad cases is good.
\end{proof}

\begin{remark}[The scope of ``forces good reduction'']
The potentially good hypothesis is essential. If $K$ contains
$\mu_{p^n}$, a split Tate curve with parameter $q=a^{p^n}$,
$0<|a|<1$, has full rational $p^n$-torsion and split multiplicative
reduction. Thus the preceding least-level assertion is not a
criterion for all elliptic curves without that hypothesis.
\end{remark}
